% !TEX program = xelatex
\documentclass[11pt,a4paper]{article}

\usepackage[UTF8]{ctex}
\usepackage[margin=2.2cm]{geometry}
\usepackage{amsmath, amssymb, amsthm}
\usepackage{listings}
\usepackage{xcolor}
\usepackage{tikz}
\usepackage{booktabs}
\usepackage{multirow}
\usepackage{hyperref}
\usetikzlibrary{positioning, arrows.meta, shapes.geometric, fit, calc, decorations.pathreplacing}

\hypersetup{
  colorlinks=true,
  linkcolor=blue!60!black,
  urlcolor=teal!70!black,
}

\lstdefinelanguage{Rust}{
  morekeywords={fn,let,mut,pub,struct,impl,use,mod,self,Self,Option,Some,None,return,if,else,for,while,loop,match,break,continue,as,where,trait,type,enum,move,ref,const,static,unsafe,extern,crate,super,dyn,async,await,box,in},
  sensitive=true,
  morecomment=[l]{//},
  morecomment=[s]{/*}{*/},
  morestring=[b]",
}

\lstset{
  language=Rust,
  basicstyle=\ttfamily\footnotesize,
  keywordstyle=\color{blue!60!black}\bfseries,
  commentstyle=\color{green!40!black}\itshape,
  stringstyle=\color{red!60!black},
  numbers=left,
  numberstyle=\tiny\color{gray},
  numbersep=8pt,
  frame=single,
  framerule=0.4pt,
  rulecolor=\color{gray!40},
  breaklines=true,
  breakatwhitespace=true,
  showstringspaces=false,
  tabsize=4,
  columns=fullflexible,
  keepspaces=true,
}

\title{\textbf{halfedge\,: 内蕴 Delaunay 三角剖分设计文档} \\ \large \texttt{src/intrinsic.rs}}
\author{halfedge 库}
\date{2026-07-05}

\begin{document}
\maketitle
\tableofcontents

\section{模块定位}

\texttt{intrinsic.rs} 实现内蕴 Delaunay 三角剖分（Intrinsic Delaunay
Triangulation, IDT），基于 Fisher, Springborn, Schröder \& Desbrun (2007)
的边翻转算法。IDT 保证所有内部边的余切拉普拉斯权重非负，
从而保证离散 Laplace-Beltrami 算子的正定性。

\subsection{与外蕴 Delaunay 的区别}

\begin{center}
\renewcommand{\arraystretch}{1.18}
\begin{tabular}{@{}lp{5.5cm}p{5.5cm}@{}}
  \toprule
  & \textbf{外蕴 Delaunay} & \textbf{内蕴 Delaunay} \\
  \midrule
  判据 & 外接圆包含对角顶点（3D 空间） & $\alpha_C + \alpha_D > \pi$（内蕴度量） \\
  几何依据 & 嵌入空间中的圆 & 曲面上的边长度量 \\
  收敛性 & 不保证 & 保证（Bobenko \& Springborn 2007） \\
  保持性 & 修改顶点位置 & 仅修改连接关系，保持内蕴度量 \\
  \bottomrule
\end{tabular}
\end{center}

\subsection{API 总览}

\begin{center}
\renewcommand{\arraystretch}{1.18}
\begin{tabular}{@{}lp{5cm}l@{}}
  \toprule
  \textbf{函数} & \textbf{功能} & \textbf{返回值} \\
  \midrule
  \texttt{intrinsic\_delaunay} & IDT 主入口 & \texttt{IntrinsicDelaunayStats} \\
  \texttt{compute\_intrinsic\_lengths} & 初始化内蕴边长 & \texttt{HashMap<HalfEdgeId, f64>} \\
  \texttt{is\_intrinsic\_delaunay\_edge} & 单边 Delaunay 判据 & \texttt{bool} \\
  \texttt{intrinsic\_cotan\_weight} & 内蕴余切权重 & \texttt{f64} \\
  \bottomrule
\end{tabular}
\end{center}

\section{核心算法}

\subsection{内蕴 Delaunay 判据}

设内部边 $e = AB$，两侧三角形为 $ABC$ 和 $ABD$，
$C$、$D$ 为对角顶点。定义 $\alpha_C$、$\alpha_D$ 为 $C$、$D$ 处的内角。

\textbf{Delaunay 条件}：
\[
  \alpha_C + \alpha_D \leq \pi
  \quad \Longleftrightarrow \quad
  \cot\alpha_C + \cot\alpha_D \geq 0
\]

当 $\cot\alpha_C + \cot\alpha_D < 0$ 时，翻转边 $AB \to CD$ 能改善
Delaunay 性质。

用边长表达：
\[
  \cot\alpha_C = \frac{l_{CA}^2 + l_{BC}^2 - l_{AB}^2}{4 \cdot \text{Area}(ABC)},
  \quad
  \cot\alpha_D = \frac{l_{AD}^2 + l_{BD}^2 - l_{AB}^2}{4 \cdot \text{Area}(ABD)}
\]

其中三角形面积由 Heron 公式：
\[
  \text{Area}(a,b,c) = \sqrt{s(s-a)(s-b)(s-c)},
  \quad s = \frac{a+b+c}{2}
\]

\subsection{翻转后新边长度}

边 $e = AB$ 被翻转为 $CD$ 后，新边的内蕴长度 $l_{CD}$ 由余弦定理计算：
\[
  l_{CD}^2 = l_{CA}^2 + l_{AD}^2
  - 2 \cdot l_{CA} \cdot l_{AD} \cdot \cos\!\left(\angle A_{ABC} + \angle A_{ABD}\right)
\]

其中：
\[
  \angle A_{ABC} = \arccos\frac{l_{AB}^2 + l_{CA}^2 - l_{BC}^2}{2 \cdot l_{AB} \cdot l_{CA}},
  \quad
  \angle A_{ABD} = \arccos\frac{l_{AB}^2 + l_{AD}^2 - l_{BD}^2}{2 \cdot l_{AB} \cdot l_{AD}}
\]

\textbf{关键性质}：翻转不改变已有边的内蕴长度，只有新边 $CD$ 的长度需要计算。
这是内蕴三角化的核心——度量只取决于曲面本身，不取决于三角化方式。

\subsection{主算法流程}

\begin{center}
\resizebox{\textwidth}{!}{%
\begin{tikzpicture}[
  node distance=0.45cm,
  proc/.style={draw, rounded corners=2pt, minimum width=11cm, minimum height=0.65cm, align=center, font=\small},
  dec/.style={diamond, draw, aspect=2.4, fill=orange!12, inner sep=1pt, font=\small, align=center, minimum width=3.5cm},
  op/.style={-Stealth, thick},
  startend/.style={draw, rounded corners=10pt, minimum width=2.6cm, minimum height=0.65cm, font=\small, fill=gray!12}
]
  \node[startend] (s) {开始：输入网格 $M$};
  \node[proc, below=of s, fill=blue!8] (p1) {1. 计算初始内蕴边长 $l_e = \|p_i - p_j\|$};
  \node[proc, below=of p1, fill=blue!8] (p2) {2. 所有内部边入队 $Q$};
  \node[proc, below=of p2, fill=blue!8] (p3) {3. 取出边 $e = AB$};
  \node[dec, below=of p3] (d1) {$\cot\alpha_C + \cot\alpha_D < 0$？};
  \node[proc, right=1.6cm of d1, fill=yellow!15] (skip) {跳过，取下一条边};
  \node[proc, below=of d1, fill=yellow!15] (p4) {4. 计算新边长度 $l_{CD}$（余弦定理）};
  \node[proc, below=of p4, fill=yellow!15] (p5) {5. 执行拓扑翻转 \texttt{flip\_edge}，更新 $l_{CD}$};
  \node[proc, below=of p5, fill=yellow!15] (p6) {6. 受影响邻接边重新入队 $Q$};
  \node[dec, below=of p6] (d2) {$Q$ 空？};
  \node[startend, below=of d2, fill=green!12] (e) {输出 Delaunay 三角化};
  \draw[op] (s) -- (p1);
  \draw[op] (p1) -- (p2);
  \draw[op] (p2) -- (p3);
  \draw[op] (p3) -- (d1);
  \draw[op] (d1) -- node[above, font=\scriptsize] {否} (skip);
  \draw[op] (d1) -- node[right, font=\scriptsize] {是} (p4);
  \draw[op] (p4) -- (p5);
  \draw[op] (p5) -- (p6);
  \draw[op] (p6) -- (d2);
  \draw[op] (d2.west) -- ++(-1.5,0) |- (p3.west)
        node[pos=0.25, left, font=\scriptsize] {否};
  \draw[op] (d2) -- node[right, font=\scriptsize] {是} (e);
  \draw[op] (skip.north) |- (d2.east);
\end{tikzpicture}
}
\end{center}

\subsection{收敛性}

Bobenko \& Springborn (2007) 证明：定义 Dirichlet 能量
\[
  E_D(\mathcal{T}) = \frac{1}{4} \sum_{e \in E(\mathcal{T})}
  \left(\cot\alpha_C(e) + \cot\alpha_D(e)\right) \cdot l_e^2
\]
每次合法翻转使 $E_D$ 严格递减，而 $E_D \geq 0$ 且可能的三角化数目有限，
故算法在有限步内终止。设安全上限为 $\text{max\_iterations} = 10 \cdot |E|$。

\section{内蕴边长系统}

\subsection{初始化}

内蕴边长初始化为 3D 欧氏距离：
\begin{lstlisting}
pub fn compute_intrinsic_lengths(mesh: &MeshStorage)
    -> HashMap<HalfEdgeId, f64>
{
    let mut lengths = HashMap::new();
    for he in mesh.halfedge_ids() {
        if let Some(l) = edge_length(mesh, he) {
            lengths.insert(he, l);
        }
    }
    lengths
}
\end{lstlisting}

twin 对的两个半边具有相同的长度。

\subsection{翻转后更新}

翻转 $AB \to CD$ 后，\texttt{flip\_edge} 保留半边 ID，
仅修改连接关系。更新 \texttt{intrinsic\_lengths}：
\begin{lstlisting}
intrinsic_lengths.insert(he, l_cd);     // he: D->C
intrinsic_lengths.insert(twin, l_cd);   // twin: C->D
\end{lstlisting}

已有边（$AC$, $CB$, $BD$, $DA$）的内蕴长度不变。

\section{复杂度}

\begin{center}
\renewcommand{\arraystretch}{1.18}
\begin{tabular}{@{}lll@{}}
  \toprule
  \textbf{阶段} & \textbf{时间复杂度} & \textbf{备注} \\
  \midrule
  初始边长计算     & $O(E)$              & 每条半边一次 \texttt{edge\_length} \\
  队列初始化       & $O(E)$              & 遍历所有无向边 \\
  单次翻转         & $O(1)$              & 余弦定理 + \texttt{flip\_edge} \\
  总翻转次数       & $O(E)$              & 实际经验值 \\
  受影响边入队     & $O(1)$ per flip     & 最多 8 条邻接边 \\
  \midrule
  \textbf{总计}   & $O(E)$（实际）      & 安全上限 $10E$ \\
  \bottomrule
\end{tabular}
\end{center}

\section{退化守卫}

\begin{itemize}
  \item 退化三角形（面积 $< 10^{-14}$）视为满足 Delaunay 条件，不翻转；
  \item \texttt{acos} 输入用 \texttt{clamp(-1.0, 1.0)} 保护，防止 NaN；
  \item \texttt{l\_cd\_sq.max(0.0).sqrt()} 防止负数开方；
  \item 最大迭代次数 $10 \cdot |E|$ 防止意外无限循环；
  \item 翻转前检查 \texttt{contains\_halfedge} 和 \texttt{is\_interior\_edge}，
        防止操作已删除或边界半边；
  \item \texttt{flip\_edge} 失败时跳过（\texttt{DegenerateTriangle} 等）。
\end{itemize}

\section{测试覆盖}

\begin{center}
\renewcommand{\arraystretch}{1.18}
\begin{tabular}{@{}ll@{}}
  \toprule
  \textbf{测试名} & \textbf{验证点} \\
  \midrule
  \texttt{intrinsic\_delaunay\_grid} & grid 翻转后所有边满足 Delaunay \\
  \texttt{intrinsic\_delaunay\_icosphere\_no\_flip} & icosphere IDT 正确 \\
  \texttt{intrinsic\_delaunay\_flat\_quad} & 扁平四边形翻转后满足 Delaunay \\
  \texttt{is\_delaunay\_equilateral} & 等边三角形满足 Delaunay \\
  \texttt{triangle\_area\_heron} & Heron 公式正确性 \\
  \texttt{flipped\_length\_diamond} & 菱形翻转后新边长度 $= 2$ \\
  \texttt{intrinsic\_lengths\_match\_euclidean} & 初始长度与欧氏距离一致 \\
  \texttt{cotan\_weight\_nonnegative\_after\_delaunay} & 翻转后余切权重 $\geq 0$ \\
  \texttt{intrinsic\_delaunay\_empty} & 空网格不 panic \\
  \texttt{intrinsic\_delaunay\_single\_triangle} & 单三角形无内部边不翻转 \\
  \bottomrule
\end{tabular}
\end{center}

全模块共 10 个单元测试；项目整体 \texttt{cargo test} 共 460 个用例。

\section{设计权衡}

\begin{itemize}
  \item \textbf{独立 HashMap vs 属性系统}：
        选择 \texttt{HashMap<HalfEdgeId, f64>} 存储内蕴边长，
        因为 IDT 算法需要频繁读写边长，属性系统的间接开销较大。
        缺点是需手动维护与网格拓扑的同步。

  \item \textbf{余切判据 vs 角度判据}：
        使用 $\cot\alpha_C + \cot\alpha_D < 0$ 而非
        $\alpha_C + \alpha_D > \pi$，因为余切可直接从边长计算，
        无需调用 \texttt{acos}（数值更稳定）。

  \item \textbf{VecDeque + HashSet 队列}：
        使用 BFS 式队列而非优先队列，因为 IDT 不需要优先级
        （任何不满足条件的边都可以翻转，顺序不影响结果）。
        HashSet 避免同一条边重复入队。

  \item \textbf{与 remesh::flip\_for\_valence 的区别}：
        remesh 翻转以优化顶点度数（外蕴目标），
        IDT 翻转以优化内蕴几何性质。两者不应混用。

  \item \textbf{安全上限 10E}：
        理论上翻转次数 $O(E)$，但设 $10E$ 安全上限防止极端情况。
\end{itemize}

\end{document}
